\section{Implementacja}
Do implementacji użyty został python z biblioteką pulp \href{https://coin-or.github.io/pulp/index.html}{https://coin-or.github.io/pulp/index.html} \\ 
Dzięki temu wykorzystujemy zarówno łatwość pythona jak i możliwości używania różnych solverów (CBC, GLPK, CPLEX, Gurobi...) przez pulpa \\ 
\begin{lstlisting}[language=Python, caption={Import bilbioteki pulp i bibliotek używanych do wykresu}]
    import pulp
    import sys
\end{lstlisting}
\begin{lstlisting}[language=Python, caption={Iniicjalizacja modelu}]
    model = pulp.LpProblem("Optimal_Distribution", pulp.LpMinimize)
\end{lstlisting}

\begin{lstlisting}[language=Python, caption={Zdefiniowanie zmiennych i parametrów}]
    fabryki = ['F1', 'F2']
    magazyny = ['M1', 'M2', 'M3', 'M4']
    dostawcy = fabryki + magazyny 
    
    klienci = ['K1', 'K2', 'K3', 'K4', 'K5', 'K6']
    
    koszt = {
        'F1': {'M1': 0.3, 'M2': 0.5, 'M3': 1.2, 'M4': 0.8,
        'K1': 1.2, 'K2': 999, K3': 1.2, 'K4': 2.0,
        'K5': 999, 'K6': 1.1},
        'F2': {'M1': 999, 'M2': 0.4, 'M3': 0.5, 'M4': 0.3,
        'K1': 1.8, 'K2': 999, 'K3': 999, 'K4': 999,
        'K5': 999, 'K6': 999},
        'M1': {'K1': 999, 'K2': 1.2, 'K3': 0.2, 'K4': 1.7,
        'K5': 999, 'K6': 2.0},
        'M2': {'K1': 1.4, 'K2': 0.3, 'K3': 1.8, 'K4': 1.3,
        'K5': 0.5, 'K6': 999},
        'M3': {'K1': 999, 'K2': 1.3, 'K3': 2.0, 'K4': 999,
        'K5': 0.3, 'K6': 1.4},
        'M4': {'K1': 999, 'K2': 999, 'K3': 0.4, 'K4': 2.0,
        'K5': 0.5, 'K6': 1.6}
    }
    P = pulp.LpVariable.dicts(
        "P", [(i, j) for i in dostawcy for j in klienci],
        cat='Binary')
    maksymalne_zamowienie = 60
    poziom_satysfakcji = {
        'K1': 50 / maksymalne_zamowienie,
        'K2': 10 / maksymalne_zamowienie,
        'K3': 40 / maksymalne_zamowienie,
        'K4': 35 / maksymalne_zamowienie,
        'K5': 60 / maksymalne_zamowienie,
        'K6': 20 / maksymalne_zamowienie}
    preferencja_klienta = {
        'K1': ['F2'], 'K2': ['M1'], 'K3': ['M2', 'M3'],
        'K4': ['F1'], 'K5': [], 'K6': ['M3', 'M4']}
    suma_satysfakcji = pulp.lpSum(
        [poziom_satysfakcji[j] * P[(i, j)]
        for i in dostawcy for j in klienci])
\end{lstlisting}

\begin{lstlisting}[language=Python, caption={Zmienne decyzyjne}]
    x = pulp.LpVariable.dicts(
        "x",
        [(i, j) for i in dostawcy for j in klienci],
        lowBound=0, cat='Integer')
    y = pulp.LpVariable.dicts("y",
    [(i, k) for i in fabryki for k in magazyny],
    lowBound=0, cat='Integer')
\end{lstlisting}

\begin{lstlisting}[language=Python, caption={Funkcje celu}]
    koszt_dystrybucji = pulp.lpSum(
        [koszt[i][j] * x[(i, j)]
        for i in dostawcy for j in klienci])
    koszt_magazynowania = pulp.lpSum(
        [koszt[i][k] * y[(i, k)]
        for i in fabryki for k in magazyny])
    alpha = 0.5
    beta = 0.5
    model += alpha
    * (koszt_dystrybucji + koszt_magazynowania)
    - beta * suma_satysfakcji
\end{lstlisting}

\begin{lstlisting}[language=Python, caption={Ograniczenia}]
for i in fabryki:
    model += pulp.lpSum(
        [x[(i, j)] for j in klienci]
        + [y[(i, k)] for k in magazyny])
        <= mozliwosci_fabryki[i]

for k in magazyny:
    model += pulp.lpSum(
        [x[(k, j)] for j in klienci]
        ) <= pojemnosc_magazynu[k]

for j in klienci:
    model += pulp.lpSum(
        [x[(i, j)] for i in dostawcy]) == zamowienia_klientow[j]
\end{lstlisting}

\begin{lstlisting}[language=Python, caption={Rozwiązanie problemu}]
    model.solve()
\end{lstlisting}
\

\begin{lstlisting}[language=Python, caption={Przedstawienie wyników}]
for v in model.variables():
    print(v.name, "=", v.varValue)
\end{lstlisting}

